% ICASSP 2027 manuscript: official SPS conference template (spconf.sty / IEEEbib.bst).
\documentclass{article}
\usepackage{spconf,amsmath,amssymb,hyperref}

\title{Beyond Reference Matching: Specification-Based Evaluation\\
of Generated DSP Implementations}

\name{Xianghui Meng$^{1}$ and Jionghao Lin$^{1,2,\ast}$}
\address{$^{1}$The University of Hong Kong \\
$^{2}$Carnegie Mellon University \\
\texttt{margretmeng1020@gmail.com} \\
$^{\ast}$Corresponding author: \texttt{jionghao@hku.hk}}

\begin{document}
\maketitle

\begin{abstract}
A reference filter is one valid solution of a design specification, not
a definition of the solution set.
On specification-defined filter-design tasks that admit more than one
valid realization, we score a candidate by specification-set
membership: $S_t=1$ if it meets a frozen passband, stopband, and
stability contract, independently of whether its coefficients match a
canonical design.
On four FIR/IIR tasks at $f_s=8\,\mathrm{kHz}$, nine generated
implementations meet the contract while failing coefficient agreement
with \texttt{firwin} or \texttt{butter} ($4/4$ tasks, $6/12$ cells).
The same contracts accept $12/12$ distinct valid controls and reject
$12/12$ mechanism mutants.
First-principles windowed-sinc and frequency-sampling designs occupy
the same set without library designers; a tighter mask still admits
two distinct same-order occupants.
Coefficient and response-to-reference oracles reject implementations
that satisfy the specification.
On these tasks, specification-set membership is the correctness
criterion we use, not agreement with one realization.
Our code is available at:
\url{https://github.com/XianghuiMeng-1020/dsp-semantic-correctness}.
\end{abstract}

\begin{keywords}
Filter design, specification testing, generated implementations, FIR, IIR
\end{keywords}

\section{Introduction}
\label{sec:intro}

DSP implementations are increasingly written by program synthesizers
and large language models~\cite{chen2021codex,liu2023evalplus}.
The default evaluation reflex is to compare the emitted coefficients,
or the emitted frequency response, with a trusted reference
implementation~\cite{virtanen2020scipy,huuhtanen2015dsp}.
That comparison is exact when the specification is a singleton: one
correct integer delay, one declared Fourier convention, one alias
period.
It is not exact for classical filter design.

A magnitude mask with a free transition band defines a
\emph{feasible set} of filters, not a unique coefficient
vector~\cite{oppenheim2010dsp,proakis2007dsp,mitra2011dsp,rabiner1975theory}.
Windowed truncation~\cite{kaiser1974window,harris1978windows},
frequency sampling~\cite{rabiner1970freqsamp},
equiripple Parks--McClellan design~\cite{parks1972chebyshev,mcclellan1973computer,rabiner1975fir,herrmann1973linear},
and IIR Butterworth or elliptic approximations~\cite{butterworth1930,constantinides1970,jackson1996filters,antoniou2006filters,vaidyanathan1993multirate}
can all meet the same passband, stopband, and stability bounds while
remaining far apart as impulse responses.
A canonical library call---Hamming \texttt{firwin}, or a Butterworth
section---is one occupant of that set.

This paper treats generated DSP code as an evaluation setting for that
structure, not as a model leaderboard.
The methodological claim is
\begin{equation}
\text{reference agreement}\;\neq\;\text{DSP correctness}
\quad\text{when}\quad
\lvert\mathcal{V}_t\rvert>1.
\label{eq:claim}
\end{equation}
We freeze a specification $S_t$, collect generated and first-principles
occupants of $\mathcal{V}_t=\{h:S_t(h)=1\}$, and compare three oracles:
coefficient concordance with one valid reference, constrained-band
response similarity to that reference, and specification membership.
When the task is specification-defined and $\lvert\mathcal{V}_t\rvert>1$,
membership is the correctness criterion used below.

\section{Specification-set evaluation}
\label{sec:method}

Let $t$ be a design task and $h$ a realized filter (FIR taps, or an
IIR pair $(b,a)$).
A frozen contract returns $S_t(h)\in\{0,1\}$.
The feasible set is
\begin{equation}
\mathcal{V}_t=\{h:S_t(h)=1\}.
\label{eq:vt}
\end{equation}
Any valid reference $h_r$ satisfies $h_r\in\mathcal{V}_t$ and nothing
stronger.
Write $\tilde h$ for the coefficient vector of $h$, or the
concatenation $(b,a)$, truncated to the shorter length when orders
differ.
Coefficient concordance with $h_r$ is
\begin{equation}
R_{t,r}(h)
=\mathbf{1}\!\left[
\frac{\lVert \tilde h-\tilde h_r\rVert_2}{\lVert\tilde h_r\rVert_2}
\le\tau_R\right],
\label{eq:r}
\end{equation}
with frozen $\tau_R=0.05$.
$R_{t,r}(h)=1$ asks whether $h$ is the same realization as $h_r$.
$S_t(h)=1$ asks whether $h$ is an admissible realization of $t$.
If $\mathcal{V}_t$ contains two realizations more than $\tau_R$ apart,
then $S_t=1$ does not imply $R_{t,r}=1$.

$S_t$ is a magnitude-mask residual on a $4096$-point \texttt{freqz}
grid at $f_s=8\,\mathrm{kHz}$.
For each constrained band $[f_0,f_1]$ with bounds $[\ell,u]$,
\begin{equation}
\rho=\max_{f\in[f_0,f_1]}
\max\bigl(\rho_\ell(f),\rho_u(f)\bigr),
\label{eq:rho}
\end{equation}
where $\rho_\ell=(\ell-\lvert H\rvert)_+/(u-\ell)$ and
$\rho_u=(\lvert H\rvert-u)_+/(u-\ell)$ with a $10^{-6}$ floor on
the denominator,
and $S_t$ uses the worst $\rho$ over bands.
FIR contracts pass at residual $\le 10^{-6}$; the IIR contract also
requires every pole radius $<0.999$ and passes at residual
$\le 10^{-3}$.
Neither floor was fit to generated residuals.
The four frozen masks are
\begin{align}
\text{LP }&
\lvert H\rvert\in[0.95,1.05]\ \mathrm{on}\ [0,800],
\nonumber\\
&\lvert H\rvert\le 0.05\ \mathrm{on}\ [2000,4000];
\label{eq:lp}\\
\text{BP }&
\lvert H\rvert\le 0.06\ \mathrm{on}\ [0,500]\cup[3200,4000],
\nonumber\\
&\lvert H\rvert\in[0.95,1.05]\ \mathrm{on}\ [1500,2200];
\label{eq:bp}\\
\text{BS }&
\lvert H\rvert\in[0.95,1.05]\ \mathrm{on}\ [0,600]\cup[3000,4000],
\nonumber\\
&\lvert H\rvert\le 0.06\ \mathrm{on}\ [1400,2200];
\label{eq:bs}\\
\text{IIR }&
\lvert H\rvert\in[0.90,1.10]\ \mathrm{on}\ [0,600],
\nonumber\\
&\lvert H\rvert\le 0.10\ \mathrm{on}\ [2400,4000],
\label{eq:iir}
\end{align}
with frequencies in hertz.
Transition bands are unconstrained, as in the classical approximation
problem~\cite{oppenheim2010dsp,rabiner1975fir}.
Linear phase is not required.

A tone battery $T$ checks consistency of the same mask on sinusoids:
real cosines of length $8192$ are passed through \texttt{lfilter}, and
steady-state Hilbert envelopes are compared with the contract bands.
$T$ is not an independent gold oracle.
Low-pass probes are $50$, $400$, $800$~(pass), $1400$~(transition),
and $2000$, $3000$, $3900\,\mathrm{Hz}$~(stop); band-pass, band-stop,
and IIR probes are the analogous edges and midpoints of
\eqref{eq:bp}--\eqref{eq:iir}.
A chirp envelope is not used as a mask residual (valid controls fail
that estimator on stopbands).
\begin{itemize}
\item \textbf{Oracle A} (reference coefficients): $R_{t,r}$.
\item \textbf{Oracle B} (response similarity): constrained-band
$\lvert H\rvert$ RMSE versus $h_r$ no larger than the maximum
same-order control--control RMSE on that task.
\item \textbf{Oracle C} (specification membership): $S_t$.
\end{itemize}
Oracle~B is a DSP-native comparison and remains a comparison to
\emph{one} element of $\mathcal{V}_t$.
Oracle~C is membership in $\mathcal{V}_t$.

\section{Experimental design}
\label{sec:design}

Generation used Qwen-2.5-72B-Instruct, Llama-3.3-70B-Instruct,
and DeepSeek-Chat at temperature $0.7$, with four draws per
task$\times$model cell.
The prompt is the contract statement.
No outcome triggered regeneration.

\noindent\textbf{RQ1 --- generated implementations.}
Canonical references are Hamming \texttt{firwin} (FIR) and
\texttt{butter} (IIR)~\cite{virtanen2020scipy}.
Before generation, three independently written valid controls per task
(\texttt{firwin}/\texttt{remez}/\texttt{firls};
\texttt{butter}/\texttt{cheby1}/\texttt{ellip}) all satisfy $S_t$.
Most same-order control pairs exceed $\tau_R$; the exceptions are
band-stop \texttt{firls}--\texttt{remez} ($\ell_2=0.005$) and IIR
\texttt{cheby1}--\texttt{ellip} ($\ell_2=0.044$).
Both exception pairs remain in $\mathcal{V}_t$.
Twelve mechanism mutants (band swap, shifted cutoff, truncated length,
unstable pole, Nyquist coded as $1\,\mathrm{Hz}$) all fail $S_t$.
The planned design is $4\times 3\times 4=48$ generations.
Same-order control pairs already realize $\lvert\mathcal{V}_t\rvert>1$:
FIR coefficient $\ell_2$ ranges are $0.14$--$0.27$ (LP),
$0.26$--$0.38$ (BP), and $0.005$--$0.051$ (BS); IIR
$(b,a)$ $\ell_2$ ranges are $0.044$--$0.611$.
Constrained-band $\lvert H\rvert$ RMSE among those pairs stays
below $0.001$ (FIR) and $0.019$ (IIR); full-grid RMSE is
$0.074$--$0.236$ (FIR) because the transition is free.
Equal-order FIR controls share group delay.

\noindent\textbf{RQ2 --- first-principles occupants.}
Two textbook FIR constructions are formed with \texttt{numpy} only.
The Type-I Hamming windowed-sinc low-pass of odd length $N$ and cutoff
$f_c$ is
\begin{equation}
h[n]=\frac{w[n]}{\sum_k w[k]}\,
2f_c\,\mathrm{sinc}\!\bigl(2f_c(n-M)\bigr),
\label{eq:sinc}
\end{equation}
$M=(N-1)/2$, $w[n]=0.54-0.46\cos(2\pi n/(N-1))$,
$f_c$ in cycles/sample~\cite{kaiser1974window,harris1978windows}.
Band-pass and band-stop designs are complementary differences of
\eqref{eq:sinc}.
Frequency sampling assigns a linear-transition desired magnitude
$A[k]$ on the $N$-point DFT grid, applies the Type-I factor
$e^{-j\pi(N-1)k/N}$, and inverts~\cite{rabiner1970freqsamp}.
SciPy design symbols are forbidden in that module.
Search is over odd $N$ and cutoff/transition placement until the
frozen residual is $\le 10^{-6}$.

\noindent\textbf{RQ3 --- tighter mask.}
One low-pass analysis mask tightens~\eqref{eq:lp} to
$\lvert H\rvert\in[0.99,1.01]$ on $[0,800]\,\mathrm{Hz}$ and
$\lvert H\rvert\le 0.01$ on $[1400,4000]\,\mathrm{Hz}$
(fivefold ripple and stop, twofold narrower free transition).
\texttt{firwin}, \eqref{eq:sinc}, and frequency sampling are compared
at a common odd length.

\noindent\textbf{RQ4 --- oracle disagreement.}
On the $14$ software-eligible generated implementations, Oracles A--C
are compared with each other and with the tone consistency check $T$.

\section{Results}
\label{sec:results}

Counts are occupancy under a frozen protocol, not prevalence.

\noindent\textbf{RQ1.}
Of $48$ requests, $20$ executed and $14$ passed the type/shape screen.
Nine of those $14$ satisfy $S_t$; all nine fail Oracle~A
(Table~\ref{tab:n}).
The event occupies $4/4$ tasks and $6/12$ generation cells;
execution outcomes are heterogeneous across cells.
Distinct occupants include a $128$-tap \texttt{firwin2} low-pass and a
$20$-tap \texttt{remez} low-pass; two $101$-tap band-stops with
relative $\ell_2>0.05$ versus \texttt{firwin}; and elliptic $(3,3)$
and Butterworth $(4,4)$ IIR designs, both stable and in-band.
Constrained tones through the nine filters are in-mask on $63/63$
probes; the twelve valid controls are in-mask on $90/90$.
On the specification bands, FIR occupants remain close in
$\lvert H\rvert$ to the canonical design (RMSE $0.0007$--$0.0045$).
Different-order FIR occupants differ in passband group delay by
$13.5$--$30.5$ samples.
The elliptic pair is not equivalent to the Butterworth reference on the mask:
constrained-band $\lvert H\rvert$ RMSE versus Butterworth is $0.065$
against a control-pair maximum of $0.019$.
Two of the nine programs have the same length as \texttt{firwin}; one
of those two still exceeds $\tau_R$ ($0.057$).
A post-hoc swap of $h_r$ among the three frozen controls leaves $7/9$
discordant with all three; the two same-length band-stops become
concordant with \texttt{remez}/\texttt{firls} while remaining
discordant with \texttt{firwin}.
The swap relocates a program to a different element of $\mathcal{V}_t$,
not out of $\mathcal{V}_t$.

\noindent\textbf{RQ2.}
Both first-principles methods enter $\mathcal{V}_t$ on all three FIR
tasks ($6/6$ shortest designs; Table~\ref{tab:fp}).
Locked to the canonical length ($81$ low-pass, $101$ band-pass and
band-stop), all six still satisfy $S_t$; five of six exceed
$\tau_R$ versus \texttt{firwin}.
Hamming windowed-sinc at $101$ taps on the band-stop is the same
method as \texttt{firwin} (coefficient $\ell_2=0.0013$,
$\Delta\mathrm{GD}=0$).
Frequency sampling at that length remains in $\mathcal{V}_t$ with
coefficient $\ell_2$ of $0.18$--$0.39$.
Every constrained tone is in-mask.
Multiplicity is a property of the mask, not of library-API diversity.

\noindent\textbf{RQ3.}
On the tighter low-pass, the shortest passing lengths are $43$
(\texttt{firwin} and windowed-sinc; numerically identical) and $57$
(frequency sampling).
At common length $N=57$, both distinct methods satisfy the tight
contract, pass $7/7$ constrained tones, share a $28$-sample group
delay, and differ by coefficient $\ell_2=0.115$ and full-grid
$\lvert H\rvert$ RMSE $0.057$, with spec-band RMSE $0.0014$
(Table~\ref{tab:fp}).
A fivefold / twofold tightening of the frozen numbers does not collapse
$\mathcal{V}_t$ to a singleton.

\noindent\textbf{RQ4.}
Table~\ref{tab:oracle} reports agreement of each oracle with the tone
consistency check $T$ (same bands as $S_t$; not an external gold).
Oracle~C and $T$ agree on all $14$ eligible implementations
($9/9$ with $S_t=1$ are fully in-mask; $5/5$ with $S_t=0$ miss at
least one constrained tone), as expected for two readings of one mask.
Oracle~A agrees with $T$ on $5/14$ and rejects every implementation
that is in-mask on $T$ ($9/9$).
Oracle~B agrees with $T$ on $8/14$ and still rejects $6/9$
in-mask implementations, namely the different-order FIR designs
and the elliptic pair.
The three filters that are in-mask on $T$ and pass Oracle~B---two
same-order band-stops and the Butterworth $(4,4)$---still fail
Oracle~A.
A reference protocol therefore accepts $0$ of $48$ planned generations
and $0/4$ tasks; the specification protocol accepts $9$ generations
and $4/4$ tasks.
On band-pass and IIR cells, ranking by coefficient distance elevates a
specification-failing implementation above a specification-valid one.
The band-pass case is a $101$-tap \texttt{firwin2} with coefficient
$\ell_2=0.401$ and residual $0.352$ (tones $5/9$) ranked above a
$128$-tap specification-valid \texttt{firwin2} with $\ell_2=1.369$
and residual $0$.
The IIR case is a specification-failing design with
$\ell_2=0.499$ versus \texttt{butter} ranked above elliptic
$(3,3)$ and Butterworth $(4,4)$ occupants with residual $0$.
Kendall $\tau$ between coefficient distance and specification residual
on the $14$ eligible rows is $0.38$.
Oracle~B calibration maxima (control-pair constrained
$\lvert H\rvert$ RMSE) are $8.6\times 10^{-4}$ (LP),
$1.0\times 10^{-3}$ (BP), $7.4\times 10^{-4}$ (BS), and
$1.9\times 10^{-2}$ (IIR).

\begin{table}[t]
\caption{RQ1 occupancy on the frozen specification-set arm.
$S_t=1$ is contract membership; $R_{t,r}=0$ is coefficient
discordance with the canonical reference.}
\label{tab:n}
\centering
\footnotesize
\setlength{\tabcolsep}{3.2pt}
\begin{tabular}{@{}lc@{}}
\hline
Quantity & Value \\
\hline
Planned / executed / eligible & $48$ / $20$ / $14$ \\
$S_t=1$ & $9$ \\
$S_t=1$ and $R_{t,r}=0$ & $9$ \\
Tasks / sources with a witness / cells & $4/4$ / $2/3$ / $6/12$ \\
Pre-generation controls ($S_t=1$) & $12/12$ \\
Pre-generation mutants ($S_t=0$) & $12/12$ \\
Constrained tones, $S_t=1$ / controls & $63/63$ / $90/90$ \\
\hline
\end{tabular}
\end{table}

\begin{table}[t]
\caption{RQ2 and RQ3.
WS: Hamming windowed-sinc.
FS: frequency sampling.
Same-order lengths are the frozen \texttt{firwin} lengths (RQ2) or
$N=57$ (RQ3).}
\label{tab:fp}
\centering
\footnotesize
\setlength{\tabcolsep}{2.6pt}
\begin{tabular}{@{}lcccc@{}}
\hline
Setting & In $\mathcal{V}_t$ & Same-order $S_t$ & $\ell_2>\tau_R$ \\
\hline
RQ2 shortest WS/FS (3 FIR tasks) & $6/6$ & --- & $6/6$ \\
RQ2 same-order WS/FS & --- & $6/6$ & $5/6$ \\
RQ3 tight LP, $N=57$, WS vs.\ FS & $2/2$ & $2/2$ & $1/1$ pair \\
\hline
\end{tabular}
\end{table}

\begin{table}[t]
\caption{RQ4: agreement with the tone consistency check $T$
($n=14$ eligible implementations).
$T$ uses the same bands as $S_t$.
$T{+}{\wedge}$oracle$-$ is oracle fail and $T$ pass;
$T{-}{\wedge}$oracle$+$ is oracle pass and $T$ fail.}
\label{tab:oracle}
\centering
\footnotesize
\setlength{\tabcolsep}{3.2pt}
\begin{tabular}{@{}lccc@{}}
\hline
Oracle & Agree with $T$ & $T{+}{\wedge}$orc.$-$ & $T{-}{\wedge}$orc.$+$ \\
\hline
A \ (coefficients vs.\ $h_r$) & $5/14$ & $1.00$ & $0$ \\
B \ (band $\lvert H\rvert$ vs.\ $h_r$) & $8/14$ & $0.67$ & $0$ \\
C \ (specification membership) & $14/14$ & $0$ & $0$ \\
\hline
\end{tabular}
\end{table}

\section{Conclusion}
\label{sec:conc}

When a specification-defined DSP task admits multiple valid
realizations, a reference implementation is one point in
$\mathcal{V}_t$.
Generated filters, first-principles windowed-sinc and
frequency-sampling designs, and a tightened magnitude mask all exhibit
the same structure: $S_t=1$ does not imply coefficient or
response agreement with a single canonical occupant.
Coefficient matching and response-to-reference matching remain useful
as realization diagnostics.
On such tasks they are not, by themselves, correctness tests.
The correctness criterion used here is specification-set membership.

\vspace{4pt}
\noindent\textbf{Compliance with ethical standards.}
No funding was received for conducting this study.
The authors have no relevant financial or nonfinancial interests
to disclose.
The study did not involve human participants, animal subjects, or
personally identifiable data.

\clearpage
\bibliographystyle{IEEEbib}
\bibliography{refs}

\end{document}
